\SourceLeaf{255}%
\Needspace{7\baselineskip}
\begin{center}
\textsc{Gauss} an \textsc{Gerling}. G\"ottingen, 23. Juni 1846.
\end{center}

\mbox{.\ .\ .\ .\ .\ .\ .} Die Unterscheidung der Schrauben h\"angt in letzter Instanz wie die Unterscheidung des Rechts und Links davon ab, dass

\begin{quote}
drei von einem Punkt $A$ ausgehende und nicht in Einer Ebene liegende gerade Linien $AB$, $AC$, $AD$ mit drei andern $ab$, $ac$, $ad$ in der Ordnung, wie sie aufgez\"ahlt sind,
\end{quote}

\begin{center}
gleiche oder verkehrte Lage haben k\"onnen.
\end{center}

\Needspace{14\baselineskip}
Nemlich im ersten Fall wird die Pyramide $ABCD$ bloss durch allm\"ahlige quantitative \"Anderungen in die Pyramide $abcd$ dergestalt \"ubergehen k\"onnen, dass $A$ in $a$, $B$ in $b$ u.\,s.\,w. \"ubergehe. Im zweiten Fall ist ein solcher \"Ubergang nicht anders m\"oglich, als dass dazwischen die ver\"anderte Pyramide einmal \textls{verschwindet} ($ABCD$ in Einer Ebene). \"Ahnliche Lagen haben

\[
\begin{array}{ccc}
AB,&AC,&AD\\
AC,&AD,&AB\\
AD,&AB,&AC,
\end{array}
\]

aber verkehrte dagegen $AB$, $AD$, $AC$ u.\,s.\,w.

Diesen \textls{Unterschied} zweier Systeme von je drei geraden Linien kann man aber \textls{nicht} auf Begriffe bringen, sondern nur aus dem Anhalten an \textls{wirklich vorhandene} r\"aumliche Dinge \textls{vorzeigen}. Zwei Geister k\"onnen sich nicht anders dar\"uber verst\"andigen, als dass sie ihre Anschauungen an Ein und dasselbe in der wirklichen Welt vorkommende System kn\"upfen. W\"ahlt man also f\"ur das 1\textsuperscript{te} System die [Ordnung:]

\Needspace{7\baselineskip}
\begin{center}
S\"ud \qquad West \qquad Zenith,
\end{center}

so ist f\"ur jedes andere Dreisystem feststehend, ob es mit diesem gleich oder ob es entgegengesetzt sei. Mit jenem gleiche so:

\Needspace{7\baselineskip}
\begin{center}
\begin{tabular}{cccc}
Vorne & Rechts & Oben & \\
Links & Vorne & Oben & \\
Vorne & Links & Unten & \\
Rechts & Vorne & Unten & \\
Rechts & Hinten & Oben & u.\,s.\,w.
\end{tabular}
\end{center}
